\section{Implications for Resource Allocation}\label{sec:implications}

These results provide a principled basis for allocating resources across the phenome. For the {{NUM_TRAITS_POWER_LIMITED}} power-limited traits, continued investment in larger GWAS cohorts is justified---each increment of sample size yields genuinely new gene discoveries. For the {{NUM_TRAITS_ANNOTATION_LIMITED}} annotation-limited traits, resources should be directed toward generating disease-relevant functional annotations: perturbation screens in relevant cell types, phenotyping in model organisms for understudied pathways, or curation of disease-specific literature. For the {{PERCENT_TRAITS_SATURATED}}\% of saturated traits, the discovery frontier has shifted from finding new associations to mechanistic understanding, and for pathway-concentrated traits specifically, indirect support provides a route to disease-relevant genes beyond the reach of GWAS alone.

% The broader point
More broadly, our analysis challenges the assumption that GWAS investment should be distributed uniformly across the phenome. The genetic discovery cycle is trait-specific, and the optimal research strategy---larger cohorts, better annotations, or mechanistic follow-up---depends on where a trait sits in this cycle. The three-level decomposition of discovery (variants, genes, mechanisms) provides a framework for making these decisions systematically rather than relying on intuition about which traits are ``well-studied'' or ``underpowered.''

% What annotation-limited means concretely
The annotation-limited category is particularly actionable. These are traits for which GWAS has identified associated loci, but current functional annotations---knockout mouse phenotypes, curated pathways, literature co-occurrence, protein interactions, and single-cell expression patterns---are insufficient to propagate signal to the full set of relevant genes. This does not mean the biology is unknowable; it means the relevant biology has not yet been encoded in a form that computational methods can leverage. Targeted experimental efforts to characterize gene function in disease-relevant contexts would directly expand the set of discoverable genes for these traits.

% Closing
We note that our saturation estimates represent a \textit{lower bound}: they are conditional on the functional annotations currently available. As new annotations are generated---particularly for the annotation-limited traits---the discovery frontier will expand, potentially reclassifying some saturated traits as having further room for discovery. The genetic discovery cycle is not fixed; it can be restarted by generating the right biological data.
